% ==============================================================================
% ACTIVIDAD 4 - FILOSOFIA DEL DERECHO / UnADM
% Producto: Estudio de caso.
% Tema: Poder, Estado y Derecho.
% ==============================================================================

\documentclass[
	spanish,
	letterpaper, oneside
]{article}

\def\documenttitle {Estudio de caso: poder, Estado y Derecho}
\def\documentsubtitle {Poder, Estado y Derecho}
\def\documentsubject {Semana 6 - Actividad 4 - Filosofía del Derecho}

\def\documentauthor {Martín Jonathan de la Cruz}
\def\coursename {Filosofía del Derecho}
\def\coursecode {FDE30}

\def\universityname {Universidad Abierta y a Distancia de México}
\def\universityfaculty {Licenciatura en Derecho}
\def\universitydepartment {Filosofía del Derecho}
\def\universitydepartmentimage {departamentos/UnADM}
\def\universitydepartmentimagecfg {height=1.57cm}
\def\universitylocation {Roma Norte, Ciudad de México}

\def\authortable {
	\begin{tabular}{ll}
		\textbf{Alumno:}
		& \begin{tabular}[t]{l}
			 Martín Jonathan de la Cruz \\
		\end{tabular} \\
		Matrícula: & ES2611202040 \\

		\textit{Profesora:}
		& \begin{tabular}[t]{l}
			Reyna Patricia \\ González Rodríguez
		\end{tabular} \\
		& \\
		\multicolumn{2}{l}{Fecha de realización: \today} \\
		\multicolumn{2}{l}{\universitylocation}
	\end{tabular}
}

\input{template}

\usepackage{helvet}
\renewcommand{\familydefault}{\sfdefault}
\usepackage{tikz}
\usetikzlibrary{calc}
\setcitestyle{authoryear,open={(},close={)}}

\def\coverwatermarkenabled {true}
\def\coverwatermarkimage {img/departamentos/UnADM.pdf}
\def\coverwatermarkopacity {0.16}
\def\coverwatermarkwidth {0.72\paperwidth}
\def\coverwatermarkxshift {0cm}
\def\coverwatermarkyshift {-0.35cm}

\newcommand{\insertcoverwatermark}{%
	\ifthenelse{\equal{\coverwatermarkenabled}{true}}{%
		\AddToShipoutPictureBG*{%
			\begin{tikzpicture}[remember picture,overlay]
				\node[
					opacity=\coverwatermarkopacity,
					inner sep=0pt,
					xshift=\coverwatermarkxshift,
					yshift=\coverwatermarkyshift
				] at (current page.center) {%
					\includegraphics[width=\coverwatermarkwidth]{\coverwatermarkimage}%
				};
			\end{tikzpicture}%
		}%
	}{}%
}

\begin{document}

\insertcoverwatermark
\templatePortrait
\templatePagecfg
\onehalfspacing

\begin{abstractd}
	\textbf{Objetivo:} Analizar la relación entre poder, Estado y Derecho mediante la definición de conceptos centrales y su aplicación al estudio de caso sobre el Estado mexicano en la Constitución Política de los Estados Unidos Mexicanos.

	\textbf{Postura de análisis:} El Estado de Derecho no se agota en que existan autoridades y leyes; exige que el poder público se origine en la soberanía popular, se distribuya institucionalmente, actúe con fundamento constitucional y encuentre límites materiales en los derechos humanos.
\end{abstractd}

\templateIndex
\templateFinalcfg

\section{Introducción}

El poder estatal sólo se vuelve jurídicamente aceptable cuando deja de presentarse como fuerza desnuda y se convierte en competencia regulada. Esta idea guía el estudio de caso: el Estado mexicano no es únicamente una organización política con territorio, población y autoridades, sino una estructura normativa que justifica, distribuye y limita el poder. Si esa estructura se rompe, el Derecho deja de funcionar como garantía y puede convertirse en simple instrumento de dominación.

La filosofía del Derecho permite observar esa tensión con más precisión. García Máynez explica que el Derecho organiza deberes, facultades y consecuencias; Kelsen ayuda a entender al Estado como una forma de orden jurídico; y Ansuátegui Roig vincula el Estado de Derecho con la protección institucional de los derechos \citep{garciaMaynez2002,kelsenTeoriaPuraDerecho2009,ansuategui2007}. En conjunto, esas perspectivas muestran que no basta preguntar quién manda, sino por qué puede mandar, con qué límites y frente a qué controles.

El caso proporcionado por la actividad se centra en los artículos 1, 39, 40, 41, 42, 49 y 133 de la Constitución. A partir de ellos se examinan soberanía, forma federal, sistema democrático, territorio, división de poderes, derechos humanos y supremacía constitucional. La posición que sostengo es que estos conceptos funcionan como una cadena: la soberanía da origen al poder, la Constitución lo organiza, la división de poderes lo controla y los derechos humanos impiden que su ejercicio desconozca la dignidad de las personas.

\section{Definición de conceptos}

\subsection{Concepto de Estado}

\textbf{Estado} es la organización política y jurídica que concentra autoridad sobre una población asentada en un territorio, mediante instituciones capaces de crear, aplicar y hacer cumplir normas. No debe entenderse sólo como gobierno, porque el gobierno cambia; el Estado permanece como estructura de organización, competencias, procedimientos y responsabilidad pública. Desde una lectura normativista, el Estado se expresa a través de un orden jurídico que permite atribuir ciertos actos a órganos públicos y no a voluntades privadas aisladas \citep{kelsenTeoriaPuraDerecho2009}. Desde una perspectiva práctica, su finalidad no puede separarse de la convivencia, la justicia y la protección de derechos \citep{lovonManualPracticoFilosofia2020}.

\subsection{Elementos del Estado}

Los \textbf{elementos del Estado} son población, territorio y poder soberano. La población integra la comunidad humana a la que se dirige el orden jurídico; el territorio delimita el espacio donde ese orden tiene vigencia; y el poder soberano permite decidir y ejecutar normas sin subordinación a otra autoridad estatal. La consecuencia jurídica de esta triada es clara: sin población no hay sujeto colectivo, sin territorio no hay ámbito espacial de validez y sin soberanía no hay autoridad capaz de producir decisiones obligatorias \citep{lovonManualPracticoFilosofia2020,cpeum2026}.

\subsection{Formas de Estado}

Las \textbf{formas de Estado} explican cómo se distribuye territorialmente el poder. Un Estado unitario concentra la decisión política en órganos centrales; un Estado federal reparte competencias entre Federación y entidades federativas; y una confederación conserva mayor independencia de las unidades asociadas. México adopta la forma federal porque el artículo 40 constitucional lo define como república representativa, democrática, laica y federal, compuesta por estados libres y soberanos en su régimen interior, unidos por el pacto constitucional \citep{cpeum2026}. Esta forma importa porque evita que toda decisión pública se concentre en un solo centro y permite reconocer ámbitos locales de gobierno.

\subsection{Sistema de Gobierno}

El \textbf{sistema de gobierno} describe la relación entre los órganos que ejercen el poder estatal. En México, el diseño constitucional es presidencial: el Poder Ejecutivo federal se deposita en una persona titular de la Presidencia, mientras el Congreso y el Poder Judicial tienen funciones propias de legislación y control constitucional. La clave filosófico-jurídica no es el nombre del sistema, sino su equilibrio: un presidencialismo sin controles se acerca al mando personal; un presidencialismo constitucional se somete a competencias, elecciones, responsabilidad y revisión judicial \citep{cpeum2026,rodriguezFilosofiaDerecho2020}.

\subsection{División de poderes}

La \textbf{división de poderes} es un principio de contención del poder. Distribuye las funciones estatales en órganos legislativos, ejecutivos y judiciales para impedir que una sola autoridad cree, ejecute e interprete la norma sin contrapesos. El artículo 49 de la Constitución mexicana expresa esa lógica al prohibir la reunión de dos o más poderes en una sola persona o corporación \citep{cpeum2026}. Su valor práctico aparece cuando una reforma aprobada por mayoría legislativa puede ser revisada por la Suprema Corte si viola el procedimiento constitucional o los derechos fundamentales.

\subsection{Estado de Derecho}

\textbf{Estado de Derecho} significa que el poder público se encuentra sometido a normas generales, procedimientos controlables, derechos exigibles y jueces independientes. No equivale a legalismo simple: una autoridad puede invocar una ley y, aun así, actuar contra el Estado de Derecho si esa ley contradice la Constitución o desconoce derechos. Ansuátegui Roig subraya que los derechos no son adorno del sistema, sino condiciones de legitimidad del poder jurídico-político \citep{ansuategui2007}. Por eso, en el sistema mexicano, el artículo 1 constitucional convierte los derechos humanos y el principio pro persona en límites materiales de toda autoridad \citep{cpeum2026}.

\subsection{Poder político, social y jurídico}

El \textbf{poder político} es la capacidad institucional de tomar decisiones obligatorias para la comunidad; el \textbf{poder social} es la influencia que ejercen actores no estatales, como medios, empresas, movimientos, grupos religiosos, sindicatos o colectivos; y el \textbf{poder jurídico} es la facultad reconocida por normas para crear, interpretar, aplicar o exigir Derecho. La diferencia es relevante: el poder social influye, el político decide y el jurídico transforma esa decisión en acto válido o inválido dentro del sistema. Ruiz Rodríguez ubica el poder como un problema central de la ontología jurídica porque conecta fuerza, legitimidad, validez y límites \citep{ruizrodriguezFilosofiaDerecho2009}. Desde esta perspectiva, el Estado de Derecho existe precisamente para que el poder político no absorba al jurídico ni convierta la influencia social dominante en imposición legal.

\section{Análisis del estudio de caso}

El estudio de caso describe al Estado mexicano a partir de preceptos constitucionales que organizan su soberanía, territorio, forma de gobierno, división de poderes, derechos humanos y supremacía constitucional. La tabla siguiente relaciona cada elemento con el artículo correspondiente y con una nota periodística que muestra cómo el concepto aparece en la práctica jurídica o política.

\begingroup
\small
\setlength{\tabcolsep}{4pt}
\renewcommand{\arraystretch}{1.35}
\begin{longtable}{@{}p{0.19\linewidth}p{0.13\linewidth}p{0.61\linewidth}@{}}
\caption{Análisis constitucional del Estado mexicano}\label{tab:analisis-caso-estado}\\
\toprule
\textbf{Elemento} & \textbf{Artículo} & \textbf{Explicación y ejemplo} \\
\midrule
\endfirsthead
\toprule
\textbf{Elemento} & \textbf{Artículo} & \textbf{Explicación y ejemplo} \\
\midrule
\endhead
\bottomrule
\endlastfoot

\textbf{Soberanía nacional} &
39 CPEUM &
La soberanía reside originariamente en el pueblo; por eso la autoridad pública no se justifica por sí misma, sino por su origen constitucional y democrático. La entrega de constancia de mayoría a Claudia Sheinbaum por el Tribunal Electoral muestra la soberanía en operación institucional: el voto popular se convierte en declaración jurídica de validez de la elección \citep{jornadaSheinbaum2024}. \\
\midrule

\textbf{Forma de Estado} &
40 CPEUM &
México se constituye como república representativa, democrática, laica y federal. El federalismo significa que las entidades tienen régimen interior propio, pero no pueden romper el pacto constitucional. La nota sobre el criterio de la Suprema Corte que reservó a la Auditoría Superior de la Federación la fiscalización de recursos federales en municipios refleja esta tensión: las competencias locales existen, pero se coordinan con límites federales \citep{infobaeAuditorias2026}. \\
\midrule

\textbf{Sistema democrático} &
41 CPEUM &
El pueblo ejerce su soberanía mediante los Poderes de la Unión y los de las entidades federativas; además, el artículo regula partidos, elecciones e instituciones electorales. La constancia presidencial de 2024 permite ver que la democracia no termina en votar: requiere cómputo, calificación, autoridad electoral y reconocimiento formal del resultado \citep{jornadaSheinbaum2024}. \\
\midrule

\textbf{Territorio} &
42 CPEUM &
El territorio nacional comprende partes integrantes como territorio continental, islas, plataforma continental, mares territoriales y espacio aéreo. La explicación periodística sobre los cruces fronterizos entre México y Guatemala permite aterrizar el concepto: la frontera no es sólo una línea geográfica, sino un espacio donde se ejerce soberanía, control migratorio, cooperación internacional y administración pública \citep{trejoCrucesFronterizos2024}. \\
\midrule

\textbf{División de poderes} &
49 CPEUM &
El poder federal se divide en Legislativo, Ejecutivo y Judicial, y esa separación impide que una sola autoridad concentre todas las funciones. La invalidación del ``plan B'' electoral por la Suprema Corte ejemplifica el control entre poderes: aunque una reforma hubiera sido aprobada por órganos políticos, el Poder Judicial podía revisar si el procedimiento legislativo respetó la Constitución \citep{elpaisPlanB2023}. \\
\midrule

\textbf{Derechos humanos como límite al poder} &
1 CPEUM &
Todas las autoridades deben promover, respetar, proteger y garantizar derechos humanos. La decisión de la Suprema Corte sobre la despenalización del aborto a nivel federal muestra que el artículo 1 funciona como límite material del poder punitivo: una mayoría legislativa no puede mantener delitos que vulneren derechos de mujeres y personas gestantes \citep{expansionAborto2023}. \\
\midrule

\textbf{Supremacía constitucional} &
133 CPEUM &
La Constitución, las leyes federales emitidas conforme a ella y los tratados son Ley Suprema de la Unión. Este principio permite resolver conflictos entre normas inferiores y la Constitución. El caso del ``plan B'' electoral también refleja esta supremacía: la Corte no sustituyó al legislador, sino que verificó si la reforma respetaba las reglas superiores que hacen válida la producción normativa \citep{elpaisPlanB2023}. \\

\end{longtable}
\endgroup

\section{Conclusiones}

Los conceptos analizados influyen directamente en la posibilidad de hablar de Estado de Derecho porque forman una estructura de legitimidad. La soberanía popular explica el origen del poder; la forma federal distribuye competencias; el territorio define el espacio de validez; el sistema democrático convierte voluntad social en decisiones institucionales; la división de poderes impide la concentración; los derechos humanos fijan límites materiales; y la supremacía constitucional ordena jerárquicamente todo el sistema. Si uno de esos elementos falla, el Estado puede conservar apariencia legal, pero pierde calidad constitucional.

La conclusión central es que el Estado de Derecho no depende sólo de tener leyes vigentes. Depende de que esas leyes nazcan de procedimientos válidos, respeten derechos, puedan ser revisadas por jueces independientes y sean aplicadas por autoridades competentes. Por eso los ejemplos periodísticos no son simples datos de actualidad: muestran que los conceptos del caso trabajan en conflictos reales, como elecciones, fronteras, fiscalización federal, reformas legislativas, control judicial y derechos reproductivos.

Desde mi postura, el criterio más importante para la práctica profesional es distinguir poder legítimo de poder meramente eficaz. Un acto de autoridad puede producir efectos y, sin embargo, ser inconstitucional si invade competencias, viola derechos o rompe el procedimiento. La formación jurídica debe servir precisamente para detectar esa diferencia. En un Estado de Derecho, obedecer la Constitución no debilita al Estado; lo vuelve jurídicamente defendible.

\bibliography{referencias/materias/filosofia-derecho/bib-filosofia-derecho/FilosofiaDerecho}

\end{document}
